\section{Guarantees and Proof Arguments}
\label{sec:guarantees}

Cross-run claims use the \emph{retry invariant}: the corrected tuple
stream, source sets, graph, scope/depth predicates, target baselines, and
resolved source/target channel admission of every outstanding unit remain
unchanged, except for baseline tags written for successful packages.
Two recurring components deserve names. \textbf{(H1)} says every baseline
channel read by \code{Propagate-Channel: inherit} remains unchanged;
\textbf{(H2)} says no contribution initially suppressed by cancellation,
hold, or correction is later released at the same $h$. H1 and H2 are
necessary in their respective cases but are not, by themselves, the full
retry invariant. Neither is needed for one invocation's determinism.

\begin{lemma}[window antitonicity]
\label{lem:w}
For $\sigma' = \mathit{run}(\sigma,F)$, $W^\circ_{\sigma'}(p)=\emptyset$
if $p$ published and $W^\circ_{\sigma'}(p)=W^\circ_\sigma(p)$ otherwise.
The train window $W$ also empties after a stable release, but may remain
unchanged after a prerelease release.
\end{lemma}
\begin{proof}
The run adds only tags at $h$. If $p$ gained none, both baseline commits
are unchanged. If it published, $\mathit{basec}_{\sigma'}(p)=h$, so
Eq.~\ref{eq:fresh} subtracts $\mathit{reach}(h)$ and is empty. Only a
stable release also makes $\mathit{scom}_{\sigma'}(p)=h$.
\end{proof}

\begin{lemma}[source stability]
\label{lem:s}
Under the retry invariant, $\mathit{src}_{\sigma_j}(u)$ is constant for
every outstanding unit throughout a valid run sequence.
\end{lemma}
\begin{proof}
The corrected tuple stream and its cancellation, hold, correction, and
source-set predicates are fixed components of the invariant. H2 is the
component excluding later release of an initially suppressed contribution.
Arbitrary tag-set inclusion outside a valid run sequence implies no such
monotonicity, because a higher-version tag may point to an earlier commit.
\end{proof}

\begin{lemma}[retry stability]
\label{lem:c}
Under the retry invariant, an unreleased package's resolved channel,
each outstanding unit's source-channel set, and
$\mathit{rsv}(u,d)$ remain constant.
\end{lemma}
\begin{proof}
These are precisely the channel and source-set components fixed by the
retry invariant. Definition~\ref{def:chan} then selects the same direct
or propagated proposal, and Eq.~\ref{eq:resolvable} reads the same sets.
\end{proof}

\begin{remarknum}
\label{rem:drift}
(H1) cannot be replaced by requiring stable completed releases: \code{inherit} reads
$\mathit{ch}(\mathit{base}_\sigma(q))$, which the origin's own
graduation moves: with \code{core} at baseline \code{1.5.0-beta.0} and
commits \code{feat(core)\^{}++1:~x} and
\code{release(core)\%stable:~graduate}, run 1 plans the consumer at
\code{2.0.1-beta.0} (inherited \code{beta}) while \code{core}
graduates; if the consumer fails, run 2 plans it at \code{2.0.1}:
same package, both plans, different versions, every completed release
stable. Thus unchanged inherited-channel state must be an explicit component of the
retry invariant; successful stable releases alone do not imply version stability.
\end{remarknum}

\begin{theorem}[the eight guarantees]
\label{thm:g}
Write $\mathit{adm}_j$ for $\mathit{adm}_{\sigma_j}$ and $n =
|\Pi(\sigma_1)|$. Under the retry invariant, for every unit $u$, consumer
$D$ and run $j \geq 1$:
\begin{description}
  \item[\G{1}] \textbf{termination.} propagation halts after $O(|\mathcal{U}|
    \cdot (|\mathcal{P}|{+}|E|))$ work, whatever the declared depths.
  \item[\G{2}] \textbf{no orphans.} $\mathit{adm}_1(u,D) \Rightarrow
    \mathit{adm}_j(u,D)$ until $D$ releases at a commit containing
    $c_u$.
  \item[\G{3}] \textbf{version stability.} if $D \in \Pi(\sigma_1)$ and $D$ is not
    released before run $k$, then $D \in \Pi(\sigma_k)$ and $\nu_k(D) =
    \nu_1(D)$.
  \item[\G{4}] \textbf{no double ledger delivery.} once $D$ releases past $c_u$,
    $\mathit{adm}_j(u,D)$ never holds again; with \G{2}, ledger discharge is
    exactly-once for every contribution whose target eventually succeeds.
    This does not make external side effects exactly-once.
  \item[\G{5}] \textbf{no widening.} $\{D : \mathit{adm}_j(u,D)\} \subseteq
    \{D : \mathit{adm}_1(u,D)\}$.
  \item[\G{6}] \textbf{convergence.} $\Pi(\sigma_{n+1}) = \emptyset$, provided
    each nonempty retry admits and publishes at least one initially
    publishable, non-held package.
  \item[\G{7}] \textbf{channel discharge.} a propagated channel contribution is
    admitted only while its commit remains in $W^\circ(D)$; baseline
    equality and transition matching independently remove no-op proposals.
  \item[\G{8}] \textbf{one-way axis independence.} bump-axis suppression alone
    does not suppress a channel change. The converse is false: resolved channels
    are an input to bump admission in Eq.~\ref{eq:resolvable}.
\end{description}
\end{theorem}

\begin{proof}
\G{1}: the traversal computing $R$ marks each package seen at most
once, so it does $O(|\mathcal{P}|{+}|E|)$ work per unit regardless of
depth or cycles; acyclicity is needed for the publish order, not here.
\G{2}: compare Eq.~\ref{eq:admission} clause by clause. Source set,
depth, graph, scope, cancellation and channel admission are fixed by the
retry invariant. The consumer's fresh window is unchanged until $D$
releases (Lemma~\ref{lem:w}); that release is the only event that can
make its admission clause false.
\G{3}: $D$ gained no tag, so its baselines, windows and
$\mathit{dir}$ are unchanged; the propagated join is unchanged by \G{2}
($\geq$) and \G{5} ($\leq$), whose proofs below use only the lemmas; so
$\mathit{eff}$ is unchanged and the channel is unchanged
(Lemma~\ref{lem:c}). Thus bump-only and channel-only membership remain;
an exact \code{Release-As} remains the same because its unit is fixed by the
invariant. In every case $\nu$ is the same function of unchanged inputs.
\G{4}: by Lemma~\ref{lem:w}, $W^\circ(D)$ empties after either a stable
or prerelease tag, killing fresh direct and propagated admission. This
also makes $\mathit{pin}_\sigma(D)$ undefined. This is discharge in the tag
ledger; an interruption after an external publish
but before its tag requires content-identity reconciliation or safe publisher repetition, not
an exactly-once claim about the external operation.
\G{5}: every clause of Eq.~\ref{eq:admission} is constant except the
window, which is antitone (Lemma~\ref{lem:w}); a conjunction of
constant and antitone predicates is antitone.
\G{6}: $\Pi(\sigma_{j+1}) \subseteq \Pi(\sigma_j)$, since a member of
the later plan was not released (else its fresh windows and pins are empty and its channel equals its new baseline) and each disjunct of
$\Pi$ is then unchanged or smaller; by the premise each nonempty retry
removes at least one initially publishable package (\G{4}), so $|\Pi|$
strictly decreases to $0$ within
$n$ runs. The bound assumes progress on every counted nonempty run; it does not bound
the number of failed attempts that make no progress.
\G{7}: the channel-axis admission test uses $c_u\in W^\circ(D)$, so
Lemma~\ref{lem:w} discharges it. While fresh, equality with the baseline
and a transition's \code{<from>} check separately reject no-ops.
\G{8}: the channel pass runs first and reads no bump-axis output; the bump
pass then reads its resolved channels through Eq.~\ref{eq:resolvable}.
\end{proof}

\begin{proposition}[tagging every admitted release]
\label{prop:tagall}
If a released package with no artefact to upload and a non-empty effective
bump were versioned but not tagged, its window would never advance
(Lemma~\ref{lem:w}), it would re-enter every plan forever, and \G{6} would
fail. Hence
Definition~\ref{def:run} tags whatever is released, and a fully successful model run has an empty subsequent plan. Held
packages are excluded from admission. This necessity is relative to a
tag-only ledger; another protocol could use other durable completion records.
\end{proposition}

Three consequences explain the operational behavior. \emph{Resumption}:
after a partial failure, $\Pi(\sigma_2)=F\cup B$ at the original versions
(\G{2}--\G{5}), provided the retry invariant holds. \emph{Idempotence}:
let $\mathit{release}(s)=\mathit{run}(s,\emptyset)$ denote a fully successful
transformation with fixed reviewed inputs and valid records. Then
\[
  \mathit{release}(s)=\mathit{release}(\mathit{release}(s)),
\]
under the retry invariant, because the first transformation discharges every
admitted obligation and the second has an empty plan (\G{4}--\G{6}). Determinism of an unchanged preview is a
separate property (Theorem~\ref{thm:det}). \emph{Scope preservation}: \G{5}
makes the original plan an upper bound on retry publication under the same
invariant. This is not an authorization guarantee against changed commands,
rewritten tags, or changed credentials.

The idempotence equation concerns recorded release state. If a publisher
succeeds but its completion record is lost, its outcome must be reconciled,
or the publisher must safely accept repetition. Arbitrary shell-side effects
are not made exactly-once by a tag ledger.


\begin{table}[t]
\centering
\caption{Running example after one downstream failure. Planned versions are
assets 1.5.0, base 0.8.3, and app 2.1.1. This is a model trace, not a timing
measurement. A no-progress attempt may be inserted without violating safety;
it is excluded from the progress premise of the convergence bound.}
\label{tab:runs}
\small
\begin{tabular}{@{}llll@{}}
\toprule
Run & Plan & Recorded success & Failure \\
\midrule
1 & assets, base, app & assets, base & app \\
2 & app 2.1.1 & app & none \\
3 & empty & none & none \\
\bottomrule
\end{tabular}
\end{table}
